Unless otherwise stated, the vertex set of
a digraph will be $\nset$ for some $n\in \N$. 

The \emph{in-degree} of a vertex $v$ in a digraph $D$ is the number of arcs
of the form $(u, v)$ in $D$; similarly, the \emph{out-degree} of $v$ is the
number of arcs of the form $(v, u)$ in $D$. A vertex $v$ in a digraph $D$ is
called a \emph{sink} if the out-degree of $v$ is $0$. A vertex is
\emph{isolated} if it has no incoming or outgoing arcs. 

If $D = (V, A)$ is a digraph, and $U$ is subset of the vertices $V$ of $D$,
then the \emph{subdigraph of $D$ induced by $U$} is the digraph with vertices
$U$ and arcs $A \cap (U\times U)$. In general, a \emph{subdigraph} of $D =
(V,A)$ is any digraph $D' = (V', A')$ with $V' \subseteq V$ and $A' \subseteq
A\cap (V' \times V')$.